% ============================================================================
% sections/sec1_introduction.tex
% ============================================================================
\section{Introduction}
\label{sec:intro}

\subsection{Challenges}
\label{sec:challenges}

Retrieval-Augmented Generation systems have become the dominant
architecture for grounding large language models in external, verifiable
knowledge: a retriever pulls the top-$k$ most relevant passages from a
corpus for a given query, and a generator conditions its answer on that
retrieved context. This design reduces hallucination and lets a model
answer questions about private or fast-changing information it never saw
during training. However, the retrieval component introduces an attack
surface that does not exist in a standalone LLM deployment, and this
project's literature survey (\Cref{sec:literature}) identifies five
distinct, paper-documented forms it takes.

\begin{description}[leftmargin=0pt,style=nextline,font=\bfseries\color{navy}]
\item[Corpus poisoning attacks.] An adversary who can insert documents
  into the retrieval corpus --- through user-submitted content, a
  compromised data pipeline, or a supply-chain attack on a shared
  knowledge base --- can plant false claims that the retriever will
  faithfully surface as if they were genuine evidence. The literature on
  this family of attacks (PoisonedRAG and its successors) shows that as
  few as five well-placed documents can flip an undefended system's
  answer with over 90\% success.

\item[Compound query-and-corpus attacks (PIDP).] The PIDP-Attack combines
  a runtime query prompt-injection with universal corpus poisoning in a
  single attack: a malicious suffix appended to the user's query steers
  retrieval toward attacker-planted ``attractor'' documents in the
  corpus, without needing prior knowledge of the victim's exact query
  text. A defense that only screens the corpus, or only screens the
  query, is structurally blind to one half of this attack.

\item[Collusion and coordinated multi-document poisoning.] Multiple
  poison documents can corroborate the same false claim. RAGuard's own
  paper states this plainly as a limitation of single-document
  leave-one-out (LOO) counterfactual verification: if document A and
  document B both support the same wrong answer, removing either one
  alone leaves the other intact, the answer does not change, and neither
  document is flagged.

\item[Semantically fluent, stealth poisoning (SilentRetrieval).]
  SilentRetrieval introduces a two-stage construction --- Coordinated
  Beam Search (CBS) and Context-Adaptive Trigger Generation (CATG) ---
  that produces poison text which is linguistically fluent, statistically
  close to clean content in embedding space, and deliberately tailored to
  the anticipated query. Such text is close to indistinguishable from
  genuine content to a perplexity filter, a lexical screen, or an
  embedding-outlier detector.

\item[False consensus in layered defenses (TriShieldRAG).] TriShieldRAG's
  own reported negative result is the most important one for this
  project's design: even a three-ring, defense-in-depth pipeline
  collapses once an attack targets the shared retrieved context every
  ring inspects --- cross-model agreement can remain high while attack
  success also remains high, because every model in the panel is voting
  on the same already-poisoned evidence. Agreement is not the same thing
  as correctness.
\end{description}

\subsubsection*{The gap: no unified, honestly-evaluated defense}

Each of the six surveyed papers addresses a subset of these threats, but
their defenses operate in isolation, adopt different assumptions about
what the attacker can and cannot do, and are evaluated on different
corpora, different attack constructions, and different metrics. No single
framework in the surveyed literature composes complementary signals
across the full RAG pipeline --- query, corpus, retrieval, and generation
--- under one honest, reproducible evaluation harness that reports both
successes and residual failure rates. Closing that specific gap is this
project's stated objective (\Cref{sec:objectives}).

\subsection{Motivation for the Work}
\label{sec:motivation}

The motivation for undertaking this project is three-fold.

\begin{enumerate}
\item \textbf{Close the documented gaps.} Each prior defense has a
  paper-stated blind spot, and this project's architecture is built so
  the weakness of one mechanism becomes the input signal for the next:
  DRS detects geometric anomalies at ingestion but has no query-time
  protection, so Ring~0 adds query screening ahead of it. ShieldRAG's
  push-pull reshaping assumes benign content holds the majority, which
  collusion defeats, so Ring~2 adds an answer-contention signal that does
  not assume majority correctness. RAGuard's leave-one-out fails against
  collusion, so Ring~3 generalises it to leave-group-out. TriShield's
  rings can collapse together under shared context, so a persistent,
  cross-query Dynamic Trust Store adds a memory that a single query does
  not have.

\item \textbf{Honest, reproducible evaluation.} An early version of this
  project produced a suspiciously perfect ``0\% attack success rate''
  headline by over-blocking, which dropped clean-query accuracy to only
  20\% --- a defense that refuses almost everything is not a usable
  result, honest or otherwise. Rebuilding the evaluation around measured
  clean-data statistics, held-out calibration splits, multi-seed
  confidence intervals, and an explicit per-ring ablation replaced that
  artifact with genuinely earned numbers, including a final
  stealth-collusion attack success rate of 0.1\% rather than an
  unexamined 0.0\%.

\item \textbf{Discover genuinely new insights, not only implement cited
  ones.} The dual-signal risk router (embedding cohesion together with
  answer-vote contention) was not part of the original design; it
  emerged directly from measuring that cohesion-only routing is
  structurally blind to an attack optimised on the answer axis rather
  than the geometry axis --- precisely the failure mode TriShieldRAG's
  own paper documents as ``false consensus.'' This architectural insight
  does not depend on the TF-IDF embedding space used for evaluation here
  and should transfer to a dense neural retriever as well.
\end{enumerate}
